\documentclass[12pt]{article}
\usepackage[utf8]{inputenc}
\usepackage[T1]{fontenc}
\usepackage{hyperref}
\title{East Danish Language Resources for K-12 Curriculum Development}
\author{Compiled by Claude Sonnet 5.5 (claude-sonnet-5-5) for Andr\'as Kornai}
\date{}
\begin{document}
\maketitle

\section{Introduction}
ISO 639-3 scy (Glottocode scan1238) is Scanian (Skånska, Skånemål), the dialects of Skåne in southern Sweden. They descend from East Danish, and Skåne became Swedish only in 1658. Classification is contested: some treat Scanian as a South Swedish dialect, some as Danish, and SIL lists it as a separate language. In practice, Scanian speakers are educated in Standard Swedish. I did not find any Scanian-medium school programme.

\section{Primary Educational Resources}
No primers, textbooks or curricula in Scanian were found.

\section{Linguistic and Reference Materials}
\begin{itemize}
\item Wikipedia, Scanian (overview of history, classification and features): \url{https://en.wikipedia.org/wiki/Scanian}
\item Carl-Erik Lundbladh, \emph{Sk{\aa}nska dialektord} (2012, Institutet f{\"o}r spr{\aa}k och folkminnen, Lund), reviewed in LexicoNordica: \url{https://tidsskrift.dk/lexn/article/download/24452/21410/56570}
\item \emph{Sk{\aa}nsk-svensk-dansk ordbok} (Sk{\aa}nska Akademien, 1995, about 2,700 entries), mentioned in the sources found; I did not see the book itself.
\end{itemize}

\section{Religious and Community Resources}
The Sk{\aa}nska Akademien and Sk{\aa}nes Hembygdsf{\"o}rbund are mentioned as regional cultural bodies. I did not find a website for either in the search results, so no URL is given.

\section{Digital and Audio Resources}
No dedicated digital or audio resources were found in the searches.

\section{Curriculum Development Context}
Teaching Scanian as a language would have to build on the existing dictionaries and on the regional identity movement. Whether it should be treated as East Danish, as Swedish dialect material or as a separate language is itself contested. Further research is needed at Lund University dialectology, which was suggested by the sources.
\end{document}